% Proposed replacement for the full report's Section 2, to align it with the
% combined notebook and sections 3--10. No new section numbering is introduced.
\section{Preprocessing and evaluation protocol}
\label{sec:shared-protocol}

We use the original student data and retain the 3,630 records with resolved
outcomes: \texttt{Dropout}=1 and \texttt{Graduate}=0. The 794 \texttt{Enrolled}
records are not treated as known graduates or used to assess predictive accuracy.
\texttt{Debtor} and \texttt{Tuition fees up to date} are excluded because their
recording times cannot be verified. No missing values or exact duplicate rows
were found. Identical retained predictor profiles are kept together, including
one pair with different outcome labels, without deleting either record.

A seeded, stratified group split reserves 726 students for testing and leaves
2,904 for development. The same five stratified group folds are used for model
selection. Grouping prevents identical retained profiles from crossing a split;
stratification approximately preserves the outcome proportions. The split seed
is 42. Every learned transformation is fitted on the relevant training rows,
then applied to validation or test rows.

Nominal codes are one-hot encoded, numerical measurements are standardized,
and binary flags remain 0/1. Development comparisons selected rare-category
pooling at a training frequency below 1\% and no log transformation. Thus the
final workflow does not apply all Assignment 1 transformations unchanged.
Constant columns are removed. Engineered features and candidate feature
selection/reduction are compared within development folds. The DNN instead uses
preprocessed original retained features. Internal neural stopping groups are
described in Section 9.

Mean development average precision (AP) selects settings within each model
family. The LDA recommendation also considers development contact-budget
performance and explainability, as described in Section 8. Each alert threshold
maximizes F2 on development cross-validation predictions and is fixed before
that model's test scoring. The top-20\% contact list is a separate budget rule.
Development reuse means these selected CV scores are not nested-CV
generalization estimates. The test split was revisited during model development,
so the comparisons are exploratory and require independent validation.
